\section{Uso de Inteligencia Artificial}
\subsection{Desarrollo}
Durante el desarrollo de la aplicación, se ha hecho uso de herramientas de inteligencia artificial para asistir en 
diversas tareas, principalmente en la generación de código y documentación. 
Estas herramientas han facilitado la creación de fragmentos de código, la redacción de textos explicativos
 y la estructuración de documentos técnicos. Asímismo, consideramos que la aplicación no hubiese sido posible  
sin la ayuda de estas herramientas, dada la complejidad del proyecto y el tiempo limitado disponible.

Es importante destacar que, aunque las herramientas de inteligencia artificial han sido de gran ayuda,
 la supervisión y revisión humana han sido esenciales para garantizar la calidad y precisión del trabajo final.

En resumen, el uso de inteligencia artificial ha sido un complemento valioso en el proceso de desarrollo
 de la aplicación, permitiendo optimizar el tiempo y mejorar la eficiencia en diversas tareas.
\subsection{Documentación}
Para la redacción de la documentación del proyecto, se ha utilizado inteligencia artificial para generar
borradores iniciales de secciones específicas, como descripciones técnicas y explicaciones de funcionalidades.
Estas herramientas han ayudado a estructurar el contenido y proporcionar un punto de partida para la elaboración
de textos más detallados y personalizados.

Sin embargo, todo el contenido generado por inteligencia artificial ha sido cuidadosamente revisado y editado
para asegurar la coherencia, precisión y adecuación al contexto del proyecto.
\subsection{Presentación}
En la preparación de la presentación del proyecto, se ha empleado inteligencia artificial para diseñar
 diapositivas y sugerir formatos visuales que mejoren la comunicación de ideas clave. 
 Estas herramientas han facilitado la creación de gráficos, esquemas y otros elementos visuales que complementan
 la información presentada. Además, esta era la primera vez que utilizábamos LaTeX para realizar una presentación,
 por lo que la ayuda de la inteligencia artificial fue crucial para aprender a usar este sistema de manera
 eficiente y aprovechar sus capacidades al máximo.